\subsection{\texorpdfstring{第~\ref{sec:immersions-embeddings}~节}{浸入与嵌入}}\label{subsec:solutions-immersions-and-embeddings}

\begin{description}[font=\normalfont,leftmargin=2em,labelwidth=1.5em]
\RGsolution{ex:coordinate-submanifold-characterization}{coordinate-submanifold-characterization}
若 $M$ 是 $N$ 的 $m$ 维子流形, 则包含映射 $i:M\hookrightarrow N$ 是嵌入. 特别地, $i$ 是浸入, 因而由局部浸入定理, 对每个 $p\in W$, 存在 $M$ 中 $p$ 附近和 $N$ 中 $f(p)$ 附近的参数化 $(U,\varphi)$ 和 $(V,\psi)$, 使得
\begin{align*}
  \widehat i(x^1,\ldots,x^m)&:=(\psi^{-1}\circ i\circ\varphi)(x^1,\ldots,x^m)\\
  &=(x^1,\ldots,x^m,0,\ldots,0)\in\mathbb{R}^n.
\end{align*}
取 $\psi((U\times\mathbb{R}^{n-m})\cap V)$ 中的开集 $W\subset N$ 及由 $x=\psi^{-1}$ 给出的坐标系 $x:W\to\mathbb{R}^n$, 则
\[
  M\cap W=\{q\in W\mid x^{m+1}(q)=\cdots=x^n(q)=0\}.
\]
反之, 若对每个 $p\in M$, 都存在 $N$ 中 $p$ 附近的坐标系 $x_p:W_p\to\mathbb{R}^n$, 使得
\[
  M\cap W_p=\{q\in W_p\mid x_p^{m+1}(q)=\cdots=x_p^n(q)=0\},
\]
则取到前 $m$ 个分量的标准投影 $\pi$, 映射
\[
  \widetilde x_p:M\cap W_p\to\mathbb{R}^m:=\pi\circ x_p
\]
是赋予子空间拓扑的 $M$ 在 $p$ 附近的一个坐标系. 事实上, $\widetilde x_p$ 是到其像的同胚:
\begin{enumerate}[label=(\roman*),leftmargin=*]
\item 若 $A$ 是 $\widetilde x_p(M\cap W_p)$ 的开子集, 则 $(A\times\mathbb{R}^{n-m})\cap x_p(W_p)$ 是 $x_p(W_p)$ 的开子集, 故 $\widetilde x_p^{-1}(A)=x_p^{-1}((A\times\mathbb{R}^{n-m})\cap x_p(W_p))\cap M$ 是 $M$ 中的开集 (相对于子空间拓扑).
\item 若 $B$ 是 $M\cap W_p$ 的开子集, 则存在 $N$ 中的开集 $B'\subset W_p$, 使得 $B=B'\cap M$, 从而
\[
  \widetilde x_p(B)=\pi\left(x_p(B')\cap\{x\in\mathbb{R}^n\mid x^{m+1}=\cdots=x^n=0\}\right)
\]
是 $\mathbb{R}^m$ 中的开集.
\end{enumerate}
此外, $\{(\widetilde x_p(W_p),\widetilde x_p^{-1})\}_{p\in M}$ 构成 $M$ 的图册: 例如, 若 $L:=W_p\cap W_q\cap M\ne\varnothing$, 而 $f_{p,q}:=x_q\circ x_p^{-1}:x_p(W_p\cap W_q)\to x_q(W_p\cap W_q)$ 是 $N$ 上相应的过渡映射, 则 $M$ 上的过渡映射
\[
  \widetilde x_q\circ\widetilde x_p^{-1}:\widetilde x_p(L)\to\widetilde x_q(L)
\]
由
\[
  (\widetilde x_q\circ\widetilde x_p^{-1})(x^1,\ldots,x^m)
  =\pi(f_{p,q}(x^1,\ldots,x^m,0,\ldots,0))
\]
给出. 从坐标的选取可以清楚地看出, 包含映射是嵌入.

\RGsolution{ex:sphere-as-regular-level-set}{sphere-as-regular-level-set}
考虑映射 $f:\mathbb{R}^{n+1}\to\mathbb{R}$, 由
\[
  f(x^1,\ldots,x^n)=(x^1)^2+\cdots+(x^{n+1})^2
\]
给出. 它的导数
\[
  (\mathrm{d}f)_x=2x^1\mathrm{d}x^1+\cdots+2x^{n+1}\mathrm{d}x^{n+1}
\]
在 $x\ne0$ 时显然是单射, 因为它由非零矩阵
\[
  \bigl(2x^1\mid\cdots\mid 2x^{n+1}\bigr)
\]
表示. 因此, 1 是 $f$ 的正则值, 故 $S^n=f^{-1}(1)$ 是 $n$ 维流形 (参见 Theorem~\ref{thm:submanifold-of-regular-level-set}). 此外,
\begin{align*}
  T_xS^n&=\ker(\mathrm{d}f)_x
  =\{v\in T_x\mathbb{R}^{n+1}\mid(\mathrm{d}f)_x(v)=0\}\\
  &=\{v\in\mathbb{R}^{n+1}\mid x^1v^1+\cdots+x^{n+1}v^{n+1}=0\}\\
  &=\{v\in\mathbb{R}^{n+1}\mid\langle x,v\rangle=0\},
\end{align*}
这里使用了等同 $T_x\mathbb{R}^{n+1}\cong\mathbb{R}^{n+1}$.

\RGsolution{ex:restriction-to-submanifolds}{restriction-to-submanifolds}
设 $i:V\to M$ 是包含映射. 则 $f\circ i:V\to N$ 可微. 对每个 $p\in W$, 设 $x_p:U_p\to\mathbb{R}^n$ ($n=\dim N$) 是 $N$ 上的一个局部坐标图, 使得
\[
  W\cap U_p=\{q\in U_p\mid x_p^{k+1}(q)=\cdots=x_p^n(q)=0\},
\]
其中 $k=\dim W$ [参见 Exercise~\ref{ex:immersions-embeddings}(\ref{ex:coordinate-submanifold-characterization})]. 由 $\widetilde x_p(q):=(x_p^1(q),\ldots,x_p^k(q))$ 定义的映射 $\widetilde x_p:W\cap U_p\to\mathbb{R}^k$ 是定义 $W$ 的一个图册的局部坐标图. 因此, 对每个 $p\in W$, 映射
\[
  \widetilde x_p\circ f\circ i:(f\circ i)^{-1}(W\cap U_p)\to\mathbb{R}^k
\]
可微, 故 $f:V\to W$ 可微.
\end{description}
